\chapter{Porting in brief}
\label{app:porting}

The kernel never touches hardware. Everything it needs comes from
\code{port/<arch>/port.h}, which every kernel file includes through \code{rtos.h}. A
port is three or four files plus a make fragment: the AVR port is the smallest complete
example, the host port shows the minimum that makes the kernel run at all, and the
Xtensa port shows a CPU with register windows, a lazy FPU and a ROM loader.
\code{porting.md} in nornDocs is the full description; in outline a port provides:

\begin{enumerate}
\item \textbf{Types and constants}: the tick, event and task mask types, the stack
  sizes (\code{RTOS\_PORT\_MIN\_STACK}, \code{DEFAULT}, \code{LARGE}, \code{STACK\_WARN},
  \code{STACK\_ALIGN}), the run time counter frequency.
\item \textbf{Critical sections}: enter returns the previous interrupt state, exit
  restores it.
\item \textbf{The interrupt model}: the \code{RTOS\_ISR} macro and the interrupt tail
  that switches once at the end of an ISR; the tick interrupt.
\item \textbf{The context switch}: initialise a stack so that the first switch to it
  enters the task trampoline; save the current context and restore another.
\item \textbf{Time}: start the tick at \code{RTOS\_CFG\_TICK\_HZ}; a free running
  microsecond source for \code{rtos\_getTimeUs}; the run time counter.
\item \textbf{Robustness hooks}: halt, reset, the memory that survives a reset for the
  panic record, the reset cause, the hardware watchdog.
\item \textbf{The make fragment}: compiler, flags, the port's sources, how to upload and
  how to simulate.
\end{enumerate}

Bringing a port up is done with tc01 first (the kernel's primitives on the new CPU) and
tc10 last (its cost), and the port is accepted when the host unit tests' invariants hold
on the target through tc01 to tc06.
